\ProjectChapter{Epilogue: The Past and Future of the Languages of Engineering}{epilogue}
\index{history of programming languages}\index{future of programming}

\begin{researchquestion}
Where did the languages of engineering come from, how do they change, and what will an
engineer's programming look like when you graduate?
\end{researchquestion}

In 1957 an engineer who wanted a computer to evaluate a formula had to translate it by
hand into hundreds of machine instructions. The first Fortran compiler did that
translation automatically, and its name said so: \emph{formula translation}. Today you can
type a differential equation into a notebook and receive its exact solution, or describe a
calculation in plain English and receive a draft program in any language of these notes.
The title of these notes, \emph{From Formulas to Programs}, describes a road that
engineers have been shortening for seventy years. This epilogue looks back along that
road, draws out the patterns, and then looks ahead. The look ahead is opinion, not fact;
read it as one engineer's informed guess.

\section*{Every language began as somebody's problem}

None of the languages in these notes was designed to be a ``language of engineering''.
Each was built by people who were stuck. John Backus was tired of hand-writing machine
code for formulas. Dennis Ritchie needed an operating system that could move from one
computer to another. Cleve Moler wanted his students to use good matrix software without
first learning Fortran. Guido van Rossum wanted a language that was pleasant to read. The
creators of Julia were tired of writing every program twice, once to try the idea and
once to make it fast. The table below lists these turning points.

\begin{table}[H]
\centering\small
\caption*{\textbf{\sffamily Six turning points.} Each new language answered a problem
that the languages of its time handled badly, and each answer is still in use.}
\begin{tabularx}{\textwidth}{@{}>{\RaggedRight\arraybackslash}p{0.11\textwidth}>{\RaggedRight\arraybackslash}p{0.38\textwidth}X@{}}
\toprule
\textbf{Years} & \textbf{The problem} & \textbf{The answer}\\
\midrule
1950s & Writing machine code for formulas is slow and full of errors. & Fortran: the computer translates the formula.\\\addlinespace
1970s & Software is tied to one model of computer. & C: one source code, compiled for any machine.\\\addlinespace
1980s & Engineers want matrices and plots without writing a compiled program. & MATLAB, and later the Wolfram Language: interactive computing in the notation of mathematics.\\\addlinespace
1980s--1990s & Programs grow to millions of lines and become impossible to organise. & C++, Java and C\#: objects, libraries and safer memory.\\\addlinespace
1990s--2000s & Scientists spend more time writing code than thinking about the problem. & \py and R: readable, free, and with a library for almost everything.\\\addlinespace
2010s & Fast code and convenient code are written in different languages; memory errors cause most security failures. & Julia: speed without a second language. Rust: safety without loss of speed.\\
\bottomrule
\end{tabularx}
\end{table}

The lesson for you is practical. A language is a tool that carries the experience of the
people who built it. When a language seems strange, ask what problem its designers had.
Fortran's arrays, C's pointers, MATLAB's matrices and \py's indentation all stop looking
arbitrary once you know the answer.

\section*{Four patterns in how languages change}

\subsection*{Languages rarely die; they settle into the foundations}

Fortran is almost seventy years old and still runs weather forecasts. C is over fifty and
still runs inside nearly every device you own. New languages do not replace old ones. They
are built on top of them. When you call NumPy to solve a linear system, \py hands the work
to C, which hands it to routines first written in Fortran. The old languages have moved
down into the foundations, where few people see them and everybody depends on them. This
is why an engineer who can \emph{read} an old language is valuable: the foundations
still need inspection and repair.

\subsection*{Each generation moves the programmer closer to the problem}

Machine code describes what the processor does. Fortran and C describe a calculation.
MATLAB and NumPy describe operations on whole matrices. The Wolfram Language and SymPy
work with the formula itself. Modelica lets you write the physical equations and leaves
the solution method to the tool. At every step the programmer writes less about the
\emph{machine} and more about the \emph{problem}, and a translator of some kind fills in
the rest. The history of programming languages is largely the history of better
translators.

\subsection*{Good ideas travel from one language to the next}

Array operations started in APL in the 1960s, became popular through MATLAB, and are now
in Fortran, \py, R and Julia. Objects started in Simula as a way to simulate queues and
traffic, and are now almost everywhere. Garbage collection came from Lisp, notebooks from
Mathematica, the data frame from S and R, and \code{async}/\code{await} became widespread through C\#.
Languages borrow freely, so the major languages grow more alike with time. This is good
news for a learner: an idea you understand in \py will meet you again, with different
spelling, in the next language you learn.

\subsection*{The community decides more than the syntax does}

\py is not the fastest language in these notes, nor the most elegant. It leads because of
what surrounds it: free libraries for nearly every task, teaching material at every
level, and millions of people who can answer a question. MATLAB holds its place through
its toolboxes and Simulink, R through its statistical packages, C through fifty years of
compilers for every chip. When you choose a language, you choose its libraries, its
tools and its people. That part of a language changes slowly, which is why the popularity
tables of Chapter~\ref{ch:intro} change slowly too.

\section*{The machine keeps changing under the languages}

For about forty years, programs became faster every year without anyone changing them,
because each new processor was faster than the last. Around 2005 that stopped. Since then
computers have gained more cores, not faster ones, and much of the fastest computing has
moved to graphics processors (GPUs) that perform thousands of simple operations at once.

This shift is reshaping the languages. A loop written the 1990s way uses one core of a
machine that has dozens. Languages and libraries that can spread work over many cores and
onto GPUs have gained ground: array libraries, Julia, the machine-learning frameworks
driven from \py, and newer additions to Fortran and C++. Even \py itself is changing: its
recent versions can be built without the internal lock that for thirty years prevented
two threads from running \py code at the same time. Expect this trend to continue. The
engineer of the next decade will more often describe \emph{what} should be computed for
a whole array or a whole grid, and leave it to the tools to decide \emph{where} it runs.

A second pressure is safety. A large share of serious security failures in software come
from memory errors that C and C++ allow (Chapter~\ref{ch:c}). As software takes control of
cars, medical devices and power grids, such errors become engineering hazards. Languages
that prevent them, with Rust as the leading example for low-level work, are being
adopted for new systems, and safety standards are likely to ask for them more often. The
existing C and C++ code will not disappear, but less new code of that kind will be
written by hand.

\section*{An old argument, open again}

The largest change of your student years is the AI assistant that writes code from a
description in ordinary words (Chapter~\ref{ch:choosing}). It has reopened an argument
that is almost fifty years old, and you should know both sides of it.

\subsection*{One view: the next programming language is the one you already speak}

Andrej Karpathy, a well-known AI researcher, summed up the new situation in one line:
``The hottest new programming language is English.'' In his account, software has passed
through three stages: programs written by hand, neural networks trained on data, and now
language models steered by instructions in everyday language. Jensen Huang, the head of
the chip maker Nvidia, drew a bolder conclusion at a summit in Dubai in 2024. For years,
he said, everyone had been told that children must learn to program; in his view the goal
of computing is the opposite, to make programming unnecessary, so that ``everybody in the
world is now a programmer.'' He advised young people to put their effort into knowing a
field well: biology, manufacturing, farming, education.

There is truth in this. People who have never written a loop can now produce a working
program by describing what they want. The barrier of syntax, which stopped many
beginners, is much lower than it was.

\subsection*{The other view: precise notation is a privilege, not a burden}

In the late 1970s the Dutch computer scientist Edsger Dijkstra wrote a short note with
the title \emph{On the foolishness of ``natural language programming''}. People were
already wishing that computers could be instructed in plain English. Dijkstra argued that
the wish was mistaken. Mathematics, he pointed out, made its great advances only after it
stopped being written in words and adopted symbols. We should, he wrote, ``regard the
convenience of using them as a privilege'': with symbols, a school pupil can do
calculations that once needed a genius. Everyday language feels easy, he added, mostly
because it lets us say things that are wrong or meaningless without noticing.

This matters to engineers in a direct way. The title of these notes begins with the word
\emph{formulas}, and a formula is exactly what Dijkstra was defending: a short statement
that means one thing. ``The drag grows with speed'' is a sentence. $F = c v^2$ is a
formula; it tells you by how much, and it can be checked. No engineering standard, drawing
or contract is written in loose prose where a number or a symbol can be used.

\subsection*{Both sides are partly right}

The two views are less opposed than they sound. Plain language is a good way to
\emph{start}: to say roughly what is wanted and get a first draft. Precise notation is
the only way to \emph{finish}: to state exactly what the program must do and to confirm
that it does. Seen against seventy years of history, the AI assistant is the next
translator in the line, from machine code to formulas, from formulas to matrices, from
matrices to equations, and now from sentences to programs. Earlier translators did not
make programmers unnecessary. Fortran did not remove the need to understand the formula;
it removed the need to count memory addresses.

There is one difference from the earlier translators, and it is important. A compiler is
built to translate the same input the same way every time, and when it fails it usually
says so. An AI assistant can produce a different program each time it is asked, and it
can be wrong while sounding certain. That is why the checking skills of these notes
matter more now, not less.

\section*{Where the road may lead}

\subsection*{Writing the code was never the hard part}

When Huang's remark spread, many working programmers gave the same reply: typing the
program was never the difficult part of the job. The difficult part was, and is,
deciding exactly what the program should do. Which cases must it handle? In which units?
To what accuracy? What should happen when the input is wrong? An assistant that types
quickly does not answer these questions. If they are left open, it fills the gaps with
its own guesses, and it does so quickly and with confidence.

For this reason many software teams now describe their work differently. They treat the
\emph{specification}, a careful written statement of what must be built, what must not,
and how to verify the result, as the thing that lasts. The code becomes something that
can be generated again whenever needed. Some call the human's new role ``architect and
reviewer'' in place of ``typist''. These are opinions from practitioners, not settled
findings, but engineers will recognise the idea at once. A bridge is built from drawings,
standards and test protocols, not from a conversation. The drawing is the lasting
record; the concrete is poured to match it.

\subsection*{Strict languages and tests become the safety net}

If a machine writes the program, something else must catch its mistakes. Two tools do
this well. The first is a strict language. A compiler for Fortran, C++ or Rust refuses a
program in which the types do not match, and so do the type checkers for \py
(Chapter~\ref{ch:python}). Some observers expect such languages to gain ground for this
reason: a generated program that must pass a strict checker is easier to trust than one
that merely runs. The second tool is the test. A set of cases with known answers, written
by a person who understands the problem, says precisely what ``correct'' means. In the
future an engineer may write the tests and leave much of the program to the machine. The
habit you are building in ENG~101, checking a result against an exact solution or a
limiting case, is exactly the skill this requires.

\subsection*{Engineering has been here before}

Electronics offers an encouraging precedent. Fifty years ago engineers drew every
transistor of a chip by hand. Today design software places billions of them
automatically. Electrical engineering did not disappear, and students of the subject
still learn how a circuit works and how to draw one. They need that knowledge to tell the
software what to build and to judge what it produces. The work moved to a higher level;
the understanding underneath stayed necessary. Programming is likely to follow the same
path.

\subsection*{Programs that can be differentiated, and models that learn}

Engineers have always needed derivatives: for optimisation, for sensitivity analysis, for
fitting models to data. A growing family of tools can differentiate an entire program
automatically, however long it is. This \emph{automatic differentiation} is what trains
neural networks, and it is spreading into simulation codes. It allows a simulation to be
tuned against measurements, and physical laws to be combined with models learned from
data. Languages and libraries built with this in mind, such as JAX and PyTorch in \py and
much of the Julia ecosystem, are likely to matter more in engineering computation.

\subsection*{Describing the system and letting the tool write the program}

In Simulink, an engineer draws a controller and a code generator writes the C program
that runs in the car. In Modelica, an engineer writes the equations of a pump and a pipe
and the tool derives the simulation. This style, in which the engineer works with a model
of the system and software produces the program, has been growing for thirty years. With
better tools and with AI assistance it will probably grow faster. More engineers will
work with models, and fewer will type every line.

\subsection*{What will probably not change}

Predictions about technology are often wrong, so it is worth listing what has stayed the
same since 1957 and is likely to stay the same.

\begin{itemize}
\item \textbf{The algorithm comes first.} A clear statement of the steps is still the
  hard part. No language and no assistant can supply it if you cannot say what you want.
\item \textbf{Numbers in a computer are approximate.} Rounding, overflow and cancellation
  (Chapter~\ref{ch:howcoderuns}) follow from how numbers are stored, not from the language.
\item \textbf{Units and assumptions live outside the code.} The Mars Climate Orbiter was
  lost through a mismatch of units that no compiler detected. The engineer still has to
  know what each number means.
\item \textbf{An unchecked result is a guess.} This holds whether the program was written
  by you, by a colleague, by a vendor or by a machine.
\item \textbf{Responsibility stays with a person.} The engineer who signs the design
  answers for the calculation behind it.
\end{itemize}

\begin{keyidea}[Key idea: learn the ideas, and the languages will follow]
Languages come and go on a timescale of decades; the ideas beneath them last much
longer. Variables, loops, functions, arrays, types, tests and the limits of
floating-point arithmetic were needed in 1957 and are needed now. If you learn those
ideas well in \py, each new language costs you a few days of syntax, not a new
education.
\end{keyidea}

\section*{What this means for you}

You will graduate into a profession in which programming is as ordinary as drawing or
report writing, and in which much routine code is drafted by machines. Three habits will
serve you whatever the languages of that time turn out to be.

Should you still learn to program, if a machine can write the code? Consider reading and
writing. Few people become novelists, and nobody suggests that the rest should stay
illiterate. Few people become mathematicians, yet every engineer learns calculus,
because without it you cannot tell whether a result makes sense. Programming is becoming
a literacy of the same kind. You learn it in order to think clearly about calculations
and to judge the ones that are handed to you.

\textbf{Learn one language deeply.} For you that language is \py. Depth in one language
teaches the ideas that transfer to the rest, and it gives you a place to test your
understanding.

\textbf{Read widely.} Use these notes to become comfortable opening a Fortran, C or
MATLAB file and working out what it does. The engineer who can read the old solver, the
vendor's script and the assistant's draft can judge all three.

\textbf{Check everything that matters.} Keep asking the question that runs through these
notes: \emph{how do I know this answer is right?} The tools for answering it, an exact
solution, a limiting case, a conservation law, a second method, do not depend on any
language.

Backus's team set out to let engineers write formulas and leave the machine instructions
to the computer. Seventy years later the computer can take on far more of the writing.
What it cannot take on is knowing which formula describes the problem, and whether the
result can be trusted. That part of the work was the engineer's in 1957, and it still is.

\section*{Sources for the debate}
{\small
A.~Karpathy, ``Software Is Changing (Again)'', talk at the Y~Combinator AI Startup School,
2025; the remark about English dates from an earlier post of his.
J.~Huang, remarks at the World Government Summit, Dubai, 2024, as reported in the
technology press.
E.~W.~Dijkstra, \emph{On the foolishness of ``natural language programming''}, note
EWD~667, \url{https://www.cs.utexas.edu/~EWD/transcriptions/EWD06xx/EWD667.html}.
The views on specifications, reviewing and strict languages are drawn from practitioner
essays published in 2025 and 2026 and are opinions, not established results.\par}
